\section{What DFG is}\label{sec:learner}

This section describes the paper's object: DFG v1.3 as frozen in \repo{study/part10/hyb_portfolio3/dfg13.py}, run alone as \code{DFG(readout="logit").fit} with every other constructor argument at its default (\cref{sec:part11}). Parameter values below are those defaults. \Cref{tab:params} lists them, with the constants fixed in the code, and line numbers refer to the frozen file. \Cref{tab:dictionary,tab:params} are generated from the frozen file itself by \repo{paper_dfg/make_tables.py}, which also checks them against the theory notes \repo{study/part11/interp/theory/THEORY.md}.

% ------------------------------------------------------------------------------------------------------------------
\subsection{In one paragraph}\label{sec:oneline}

DFG, \emph{discrete feature growth}, is one greedy loop over an explicit dictionary of Boolean features. It starts from the constant. In each round it fits the $\pm1$ labels by least squares on the features it has, and takes the residual $r=y-\hat y$ and its root mean square $s$. It then enumerates candidate features class by class: Fourier characters $\chi_S(x)=\prod_{i\in S}x_i$, and indicators of the values of signed sums $\sum_{i\in S}s_ix_i$. A candidate is admitted only if its correlation with $r$ passes a test whose threshold is a union bound over the candidate's whole class, about $s\sqrt{2\ln N/n}$ for a class of $N$ candidates. (At $n\le512$ a forced variant may add the best candidate when nothing passes; DFG decides by cross-validation whether to use it.) DFG refits the linear readout on the enlarged set of features and repeats, until a round admits nothing or the fit has every training label's sign right. It predicts with the sign of that readout, refined by a few logistic Newton steps, plus one Hebbian step on the final residual. The fitted model is therefore a short list of named features with real weights. For example, with seed 0, DFG's model of par4 $=x_1x_2x_3x_4$ at $n=64$ is the constant and $\chi_{\{1,2,3,4\}}$. Its model of sym9 $=[\,|x_1+\dots+x_9|\le1\,]$ at $n=128$ is the constant and one folded block, three indicators of the value of $|x_1+\dots+x_9|$. Both have test error 0 (\repo{study/part11/interp/tables/atlas.md}; \cref{sec:interp}). \Cref{fig:loop} gives the loop in outline.

\begin{figure}[t]
\centering
\fbox{\begin{minipage}{\dimexpr\linewidth-2\fboxsep-2\fboxrule\relax}\small
\textbf{fit}$(X,y)$, with $n$ rows and $d$ inputs. If $n>512$, return \textbf{grow}$(X,y;\,\text{default},1)$. Otherwise run \textbf{grow} for both variants, default and forced, on each of $F$ cross-validation folds ($F=10$ at $n\le64$, else 5), with budget scale $1/(4F)$ each. Let $e_{10}$ ($e_{01}$) count the held-out points where only the default (only the forced) variant errs. Return \textbf{grow}$(X,y;\,\text{forced},\frac12)$ if $e_{10}-e_{01}>2\sqrt{e_{10}+e_{01}}$, else \textbf{grow}$(X,y;\,\text{default},\frac12)$.

\smallskip
\textbf{grow}$(X,y;\,\text{variant},\text{scale})$. Start with the features $\Phi=\{1\}$. In rounds $0,1,\dots,7$:
\begin{enumerate}[leftmargin=2.2em,label=\arabic*.,itemsep=1pt,topsep=2pt]
\item fit the weights $w$ by least squares on $\Phi$ (with a small ridge); $r=y-\Phi w$, $s=\mathrm{rms}(r)$;
\item stop if 0.75 of the budget is spent, or, from round 1 on, if $\mathrm{sign}(\Phi w)=y$ on every training row;
\item characters: in each class $\mathcal C$ (global, extension, local; \cref{tab:dictionary}), find the $\chi_S$ with $|\mathrm{mean}(r\chi_S)|>s\sqrt{2\ln(4K|\mathcal C|/\delta)/n}$;
\item dense moves: the linear block (every remaining $x_i$, by one $\chi^2$ test), and spectral features $(u\cdot x)^2-1$ for the eigenvectors $u$ of the residual pair matrix above its noise edge;
\item sum blocks: search for a signed subset $(S,s)$ on $r$; keep it if the projection of $r$ on its group indicators passes a $\chi^2$ test with a union bound over every $(S,s)$ of its size, both on $r$ and on $r$ projected off this round's characters; likewise products, modulated and weighted blocks;
\item if nothing passed: stop (default variant), or add the best candidate anyway (forced variant, at most twice per fit);
\item add what passed to $\Phi$, keeping $|\Phi|\le\min(256,\lfloor n/2\rfloor)$.
\end{enumerate}
\textbf{Readout.} If $n\ge8|\Phi|$, refine $w$ by 15 Newton steps of logistic regression. Predict $\mathrm{sign}\bigl(\Phi(x)w+g\cdot x\bigr)$, with $g=X^\top r/n$ on the final residual.
\end{minipage}}
\caption{One call of DFG v1.3, in outline (\repo{study/part10/hyb_portfolio3/dfg13.py}: \code{DFG.fit}, lines 548--571; \code{DFG._fit}, lines 580--877). Here $K=4$ and $\delta=0.1$. \Cref{sec:round,sec:smalln} give every step exactly.}
\label{fig:loop}
\end{figure}

% ------------------------------------------------------------------------------------------------------------------
\subsection{The hypothesis class}\label{sec:hypothesis}

DFG predicts $\mathrm{sign}\,h(x)$, and $-1$ where $h(x)=0$, with
\begin{equation}\label{eq:dfgclass}
h(x)=\sum_{S\in\mathcal S}w_S\,\chi_S(x)+\sum_{i\in\mathcal L}w_i\,x_i+\sum_{B\in\mathcal B}\sum_{v\in\mathcal V_B}w_{B,v}\,\chi_{P_B}(x)\,\bigl[t_B(x)\in G_{B,v}\bigr]+\sum_{u\in\mathcal U}w_u\bigl((u\cdot x)^2-1\bigr)+g\cdot x .
\end{equation}
The five sums are the five parts of the dictionary.
\begin{itemize}
\item \emph{Characters.} $\mathcal S$ holds the constant, $\chi_\emptyset=1$, and every character admitted by its own test or forced (\cref{sec:smalln}). The global classes reach degree 6. The extension and local classes, which grow from the characters already present, reach higher degrees.
\item \emph{The linear block.} $\mathcal L$ is a set of inputs whose degree-1 characters were admitted together, by one test on all of them. In the code they are ordinary degree-1 characters (line 832). They are written apart here because they entered together, without tests of their own.
\item \emph{Sum blocks.} A block $B$ has inputs $S_B$, signs $s_i\in\{\pm1\}$, and the integer statistic
\[
t_B(x)=\sum_{i\in S_B}s_ix_i+a_B\,\tau_B(x),
\]
or its absolute value if the block is \emph{folded}. An unweighted block has $a_B=0$. A \emph{weighted} block (v1.3) has $a_B\in\{\pm2,\pm3\}$ and an inner term $\tau_B$ on inputs outside $S_B$: a literal $x_j$, a 2-parity $x_jx_l$, or an AND of two literals written as $\pm1$. The groups $G_{B,v}$ are contiguous runs of the values of $t_B$ that occur in training, each holding at least $\max(3,\lfloor0.005\,n\rfloor)$ training points. Rare values are merged into their neighbours (\code{_edges}, lines 86--97). $\mathcal V_B$ is every occupied group except the most populous one, whose indicator the constant replaces. A \emph{modulated} block multiplies its indicators by a character $\chi_{P_B}$ that is already a feature, on inputs outside $S_B$; otherwise $\chi_{P_B}=1$. A \emph{product} block $B_1\times B_2$ of two blocks with disjoint inputs uses the joint cells $G_{B_1,v_1}\times G_{B_2,v_2}$ of the pair $(t_{B_1}(x),t_{B_2}(x))$. When every value of $t_B$ is its own group, every function of $t_B$ is a combination of the constant and the block's indicators (\cref{prop:blocks}). A function of a signed sum of $k$ inputs then costs at most $k+1$ weights, where its Fourier expansion can need all $2^k$ characters on $S_B$.
\item \emph{Spectral features.} Each $u\in\mathcal U$ is a unit eigenvector of the residual pair matrix (\cref{eq:spectest}). The feature $(u\cdot x)^2-1=2\sum_{i<j}u_iu_jx_ix_j$ is a fixed combination of degree-2 characters.
\item \emph{The kernel background.} $g=\frac1n\sum_tr_tx_t$ is computed on the final least-squares residual: one Hebbian step over all $d$ inputs. It is computed, not fitted.
\end{itemize}
The weights $w$ are fitted on the $p\le\min(256,\lfloor n/2\rfloor)$ admitted features, by least squares and then, when $n\ge8p$, by logistic Newton steps (\cref{sec:round}).

\begin{table}[t]
\centering{\footnotesize\setlength{\tabcolsep}{4pt}
\begin{tabular}{@{}lllrrr@{}}
\toprule
class & candidates & rows tested & $N$ at $d=30$ & $\ln(4KN/\delta)$ & \hd{smallest $n$\\that can pass}\\
\midrule
\multicolumn{6}{@{}l}{\emph{Characters} $\chi_S$, test \eqref{eq:chartest}}\\
global, $k=1$ & every $x_i$ & all & 30 & 8.48 & 16.95\\
global, $k=2$ & every pair & all & 435 & 11.15 & 22.30\\
global, $k=3$ & every triple & all & 4\,060 & 13.38 & 26.77\\
global, $k=4$ & every 4-set; round 0 & all & 27\,405 & 15.29 & 30.59\\
global, $k=5$ & round 0, final fit & first 1\,374 (687) & 142\,506 & 16.94 & 33.88\\
global, $k=6$ & round 0, final fit & first 147 (73) & 593\,775 & 18.37 & 36.74\\
extension, $+1$ input & $T\cup\{j\}$ & all & $bd\le960$ & 8.48--11.94 & 16.95--23.88\\
extension, $+2$ inputs & $T\cup\{j,l\}$ & all & $b\binom d2\le13\,920$ & 11.15--14.62 & 22.30--29.23\\
local & $S\subseteq V$, $|V|\le12$ & all & $\binom{|V|}{|S|}\le924$ & 5.08--11.90 & 10.15--23.81\\
\addlinespace
\multicolumn{6}{@{}l}{\emph{Dense moves}}\\
linear block & every remaining $x_i$ & all & \multicolumn{2}{r}{$\chi^2_d$ threshold 51.2, \eqref{eq:lintest}} & 62\\
spectral features & $\le4$ eigenvectors & all & \multicolumn{2}{r}{edge \eqref{eq:spectest}} & 120\\
\addlinespace
\multicolumn{6}{@{}l}{\emph{Sum blocks} on $k$ inputs, test \eqref{eq:blocktest}, $\mu=0$; searched on the first 2\,048 rows}\\
$k=2$ & $(S,s)$, folded or not & all & $4.2\times10^{4}$ & 15.71 & 40.4\\
$k=3$ & $(S,s)$, folded or not & all & $7.8\times10^{5}$ & 18.64 & 46.9\\
$k=5$ & $(S,s)$, folded or not & all & $1.1\times10^{8}$ & 23.59 & 57.9\\
$k=9$ & $(S,s)$, folded or not & all & $1.8\times10^{11}$ & 30.97 & 74.1\\
$k=14$ & $(S,s)$, folded or not & all & $5.7\times10^{13}$ & 36.75 & 86.6\\
$k=24$ & $(S,s)$, folded or not & all & $2.4\times10^{14}$ & 38.18 & 89.7\\
weighted & $+a\tau$, $a\in\{\pm2,\pm3\}$ & all & $\times\,17\,640$ & $+9.78$ & \\
every block of a round & $\mu\le4$ modulators $\chi_P$ & all & $\times\,(1+\mu)$ & $\le+1.61$ & \\
product & $A\times B$, disjoint & all & $N_AN_B$ & $L_A+L_B$ & \\
\bottomrule
\end{tabular}
}
\caption{DFG's dictionary at $d=30$ \status{Proved, by the code}. Each row is a class: a set of candidates enumerated together and tested with one union bound. ``Rows tested'': a budgeted search runs on the first $m$ rows, with $m$ set by its share of the budget (\cref{sec:smalln}). At $d=30$ that is every row of the grid ($n\le16384$) except for degrees 5 and 6, which see the first 1374 and 147 rows (687 and 73 under the cross-validation's half budget, $n\le512$). $N$ is the count the union bound pays for, and $\ln(4KN/\delta)$ the log term of the threshold, with $K=4$ and $\delta=0.1$. For the extensions, $b\le32$ is the number of current characters extended; for the local class, $V$ is the set of inputs of the current characters. ``Smallest $n$ that can pass'': no member can pass its test on fewer rows, for any labels (\cref{prop:floor}). For characters this is $2c^2\ln(4KN/\delta)$; for blocks it is $1+2\sqrt{L_B}+2L_B$, the threshold of \cref{eq:blocktest} at $\nu=1$. For the linear block it is the room condition $n\ge2(p+k)$, which is 62 in round 0 with all 30 inputs; for the spectral features it is $n\ge4d$. For sum blocks, $N_B$ counts the supports, the sign patterns up to a global sign, the two fold variants, the $1+\mu$ choices of modulator and the $k_{\max}=24$ sizes. Weighted and modulated blocks multiply it by the factors shown, and a product pays for both factors. Source: \repo{study/part10/hyb_portfolio3/dfg13.py} via \repo{paper_dfg/make_tables.py} (\texttt{tab\_dictionary}; values in \repo{paper_dfg/tables/numbers.json}).}
\label{tab:dictionary}
\end{table}

\paragraph{The dictionary.} \Cref{tab:dictionary} lists the classes. The dictionary is explicit and finite. At $d=30$ the global classes hold $\sum_{k\le6}\binom{30}k=768\,211$ characters. The unweighted blocks $(S,s,\text{fold})$ on 2 to 24 inputs number $2.0\times10^{14}$, close to $3^{30}=2.1\times10^{14}$, the number of signed subsets of 30 inputs. A class's size $N$ enters its test only through the log term $\ln(4KN/\delta)$. The global classes are fixed by $d$. The extension and local classes are built from the characters already present, so they are small and their thresholds are low. Features on inputs already in use are therefore cheaper to admit: the discrete counterpart of the composition effect that Parts 1--5 saw in the networks \cite{gdcircuit2026}. A block on $k$ inputs pays $\ln\binom dk+k\ln2+\ln\bigl(4K(1+\mu)k_{\max}/\delta\bigr)$. That is $k\ln2+\ln(24(1+\mu))$ more than a character of degree $k$, although the block class of size $k$ is $2^k$ times larger.

\paragraph{Which parts matter} \status{Numerical; post hoc}. Not every part is needed for DFG's results. The protocol's addendum, written after every pre-registered result had been read (\repo{study/part11/PROTOCOL.md}), defines a \code{core} run. It keeps the global characters of degree at most 3, the unweighted sum blocks (seeded by pairs and single inputs), the residual-scaled tests, the least-squares refit and the small-$n$ cross-validation. It drops the linear block and the spectral features, the extension and local classes, the degree-4 to degree-6 searches, the weighted blocks, the eigenvector seeds, the logistic refinement and the kernel background. Paired with DFG alone on the same training sets, the core's mean $\log_2\nstar$ ratio is $+0.07$ (median $-0.01$) over the 133 of 138 targets both learn. It needs clearly less data on 28 and clearly more on 26, and it loses par5, which needs the degree-5 search (\repo{study/part11/RESULTS.json}, \texttt{q5b\_core}; \cref{sec:ablations}). On these suites, then, the loop and its tests, with low-degree characters and signed-sum blocks, carry most of the results.

% ------------------------------------------------------------------------------------------------------------------
\subsection{One round}\label{sec:round}

A round has five steps: the residual, the character tests, the dense moves, the sum blocks and the refit. Each step uses only statistics of the current residual on the training rows.

\paragraph{Residual and noise level.} The features are the columns of $\Phi\in\R^{n\times p}$, the first one constant. The weights minimize $\sum_t(y_t-(\Phi w)_t)^2+10^{-3}\,n\sum_{j\ge2}w_j^2$ (line 589). The constant is not penalized, so the residual $r=y-\Phi w$ has mean 0 up to rounding. The noise level is $s=\sqrt{\max(\mathrm{mean}(r^2),10^{-12})}$.

\paragraph{The character test.} A candidate $S$ belongs to a class of $N$ candidates, tested on the first $m$ rows ($m=n$ unless the class is subsampled). It passes if
\begin{equation}\label{eq:chartest}
z_S=\frac{|v_S|}{c\,\sigma_m\sqrt{2L_N/m}}>1,\qquad v_S=\frac1m\sum_{t\le m}r_t\,\chi_S(x_t),\qquad L_N=\ln\frac{4KN}{\delta},
\end{equation}
with $\sigma_m=\mathrm{rms}(r_1,\dots,r_m)$ (so $\sigma_n=s$), $c=1$, $\delta=0.1$ and $K=4$ in every round (lines 602--603). $K_{\rm later}=3$ caps only the global degree after round 0. The standard error of $v_S$ is about $\sigma_m/\sqrt m$, so the test asks for $c\sqrt{2L_N}$ standard errors: a Bonferroni threshold over the class's $N$ candidates. The classes are:
\begin{itemize}
\item \emph{global}: every $\chi_S$ with $|S|=k$, for $k\le K=4$ in round 0 and $k\le K_{\rm later}=3$ after; $N=\binom dk$. Degrees 1 and 2 come from the degree-1 coefficients $g=X^\top r/n$ and the residual pair matrix $M=\frac1n\sum_tr_tx_tx_t^\top$. Degrees 3 and 4 come from the same matrix with the rows weighted by $x_a$ or $x_ax_b$ (lines 616--648). In round 0 of the final fit, degrees 5 and 6 ($K_{\rm hi}=6$) come from $\frac1m\sum_tr_t\chi_{T_1}(x_t)\chi_{T_2}(x_t)$ over disjoint $T_1,T_2$ of degrees $\lceil k/2\rceil$ and $\lfloor k/2\rfloor$, one product of two character tables (lines 651--663);
\item \emph{extension}: for each of the up to 32 current characters $\chi_T$ with the largest $|w_T|$, the sets $T\cup\{j\}$ ($N=bd$ for $b$ such bases) and $T\cup\{j,l\}$ ($N=b\binom d2$), from the degree-1 and pair statistics of $r\chi_T$ (lines 665--681);
\item \emph{local}: if the current characters use $2\le|V|\le12$ inputs, every $S\subseteq V$, from one Walsh--Hadamard transform of $r$ on $V$; $N=\binom{|V|}{|S|}$ (lines 683--692).
\end{itemize}
A candidate that passes in several classes keeps its largest $z_S$. The round admits the up to 64 candidates with the largest $z_S$, within the room $\min(256,\lfloor n/2\rfloor)-p$.

\paragraph{Dense moves.} Two moves catch signal spread over many small coefficients, none of which would pass a test of its own (lines 694--709). Let $\mathrm{lin}$ be the inputs whose character is neither a feature nor a candidate this round, and $k=|\mathrm{lin}|$. The \emph{linear block} admits all of them at once if
\begin{equation}\label{eq:lintest}
T_{\rm lin}=\frac n{s^2}\sum_{i\in\mathrm{lin}}g_i^2>k+2\sqrt{k\ln(1/\delta)}+2\ln(1/\delta)\quad\text{and}\quad n\ge2(p+k).
\end{equation}
This is Laurent and Massart's tail bound for $\chi^2_k$ at level $\delta$ \cite{laurent2000adaptive}. The block carries halfspaces and majorities. The \emph{spectral features} take the eigenvectors $u$ of $\tilde M$, the symmetrized $M$ with zero diagonal and with the entries of pairs that are features or candidates set to 0, whose eigenvalues satisfy
\begin{equation}\label{eq:spectest}
|\lambda|>2s\sqrt{d/n}\,\Bigl(1+\sqrt{2\ln(1/\delta)/d}\Bigr),
\end{equation}
at most four of them, and only if $n\ge4d$. The first factor is the edge of the semicircle law for a $d\times d$ Wigner matrix with entries of size $s/\sqrt n$, about the size of $\tilde M$'s entries under noise. Each such $u$ adds the feature $(u\cdot x)^2-1$.

\paragraph{The sum-block test.} Take a block $B$ with occupied groups $v$, each of $n_v$ training points, and a vector $\rho$: the residual $r$, or $r\chi_P$ for a block modulated by $\chi_P$. The block's statistic is the squared norm of the projection $\Pi_B\rho$ of $\rho$ on the block's group indicators, in units of $s^2$. The block passes if
\begin{equation}\label{eq:blocktest}
T_B(\rho)=\frac1{s^2}\sum_vn_v\,\bar\rho_v^{\,2}=\frac{|\Pi_B\rho|^2}{s^2}\;>\;\nu+2\sqrt{\nu L_B}+2L_B,
\end{equation}
\[
L_B=\ln\frac{4KN_B}\delta,\qquad N_B=\tbinom dk\,2^{k-1}\cdot2\cdot(1+\mu)\cdot k_{\max},
\]
where $\bar\rho_v$ is the mean of $\rho$ over group $v$, $\nu=\max(\text{occupied groups}-1,1)$, $k=|S_B|$, $\mu\le4$ is the number of current characters used as modulators, and $k_{\max}=24$ (line 719). $N_B$ counts the supports, the $2^{k-1}$ sign patterns up to a global sign, the two fold variants, the $1+\mu$ modulators and the 24 sizes. By Laurent and Massart, $\Pr(\chi^2_\nu\ge\nu+2\sqrt{\nu L}+2L)\le e^{-L}$. With Gaussian noise, a fixed candidate would therefore pass with probability at most $e^{-L_B}$. The union over every $(S,s,\text{fold},\text{modulator},\text{size})$ would then be at most $\delta/(4K)$, whichever candidates the search visits. That calibration is Gaussian; it is not proved for DFG's residuals, whose groups depend on the data (\cref{sec:guarantees}). A weighted block adds $\ln\bigl(2(4d+20\binom d2)\bigr)=9.78$ nats at $d=30$ for its inner term (line 136). The count is: $a\in\{\pm2,\pm3\}$ times a literal on any input, or a 2-parity or an AND of two literals on any pair, times 2 for a separate share of $\delta$. A product pays $L_{B_1}+L_{B_2}$. A block is searched on the round's residual, but admitted only if it also passes the same threshold on $\tilde r$, the residual projected off the constant and this round's passing characters (up to 64), with $\mathrm{mean}(\tilde r^2)$ in place of $s^2$ (lines 725--738). So it must add something beyond those characters. If $\mathrm{mean}(\tilde r^2)\le10^{-4}s^2$, the characters already explain the residual, and no block is searched.

\paragraph{The block search.} The test does not say which $(S,s)$ to test. \code{_block_search} (lines 179--500) proposes candidates from the residual alone. It scores them by their margin, the statistic minus its threshold, on the first $\min(n,2048)$ rows. Its starting points are:
\begin{itemize}
\item prefixes of the inputs ordered by $|u_i|$, for the top two eigenvectors $u$ of $\tilde M$, with signs $\mathrm{sign}(u_i)$; every prefix length $k\ge2$, folded and unfolded. If $f$ is a function of $t=\sum_{i\in S}s_ix_i$, then $\E[f\,x_ix_j]=a_2s_is_j$ for distinct $i,j\in S$ and 0 for every other pair (\cref{prop:blocks}), so $\tilde M\approx a_2(ss^\top-I)$ on $S$, and its top eigenvector gives both the support and the signs;
\item prefixes of the inputs ordered by $|g_i|$, with signs $\mathrm{sign}(g_i)$;
\item greedy growth from the 3 strongest pairs of $\tilde M$ (both relative signs) and the 2 strongest single inputs. Each step adds the input and sign that maximize $T_B$, until 24 inputs or 5 steps without a better margin;
\item at $n\le128$, in a round's first additive search, a beam of width 32 over $(S,s)$, started from the best pairs and scored on the first $\min(n,1024)$ rows;
\item in the additive search, v1.3's weighted candidates. These are eigenvector prefixes whose first input has weight 2 or 3; an inner term $\tau\in\{x_jx_l,\mathrm{AND}(\pm x_j,\pm x_l)\}$ on the strongest pair of $\tilde M$, with $a\in\{\pm2,\pm3\}$; and extensions $t+a\tau$ of the two best unweighted candidates. The extensions use a literal outside $S$, an input of $S$ re-weighted, or a 2-parity or AND on one of the 24 pairs outside $S$ with the largest $|\tilde M_{jl}|$.
\end{itemize}
The three best candidates are refined by single moves (add an input with either sign, remove one, flip a sign) while the margin improves, at most 72 moves. Weighted candidates are refined in the same way, with their inner term fixed. The best weighted candidate replaces the best unweighted one if its margin is larger. The winner's groups are then formed on all $n$ rows, and the winner is tested by \cref{eq:blocktest} on all $n$ rows.

\paragraph{Block moves in a round.} At $n\ge32$ (lines 715--815):
\begin{itemize}
\item up to two \emph{additive} blocks. The second is searched after the first candidate's group indicators are projected out of the residual, whether or not the first passed;
\item \emph{products}. When the round has two additive candidates, a block $B$ is searched over the joint cells with the first one, $A$, starting also from the second. When $A$ has at least two inputs of each sign, a \emph{split} is also tried: $B$ is searched over the joint cells with the block of $A$'s positive inputs, starting from its negative inputs. The better product is admitted only if its margin on $\tilde r$ exceeds $A$'s. Otherwise the product of the round's two additive candidates is tested, if their inputs are disjoint;
\item one \emph{modulated} block for each of the up to 4 current characters $\chi_P$ with the largest $|w_P|$, searched on $r\chi_P$ with $P$'s inputs excluded;
\item the products of pairs of additive blocks, admitted in earlier rounds, that have disjoint inputs.
\end{itemize}

\paragraph{Refit and stopping.} The round adds the admitted characters, the linear block's inputs, each admitted block's indicators (if they fit in the room) and the spectral features, and refits the weights. DFG stops (lines 596--598, 824):
\begin{itemize}
\item when a round admits nothing (in the default variant; \cref{sec:smalln} for the forced one);
\item when no room is left, $p\ge\min(256,\lfloor n/2\rfloor)$;
\item from round 1 on, when the fit has every training label's sign right;
\item at the start of a round, if more than 0.75 of the call's budget is spent;
\item after 8 rounds.
\end{itemize}
The code also stops if $s<10^{-6}$. That cannot happen, since $s\ge10^{-6}$ by its definition (\repo{study/part11/interp/calib/VERIFY_CALIB.md}).

\paragraph{Readout.} After the last round, if $n\ge8p$, the least-squares weights are refined by 15 Newton steps on the mean logistic loss plus $\frac12\cdot10^{-4}|w|^2$, started from $4w$. The result is rescaled so that its output on the training rows has the norm of the least-squares output (lines 849--855). Part 6 added this step after the least-squares readout plateaued at 5--7\% test error on a dense halfspace (thr30) where the networks reached 1\% \cite{discretelearners2026}. The prediction adds the kernel background $g\cdot x$ (line 856).

% ------------------------------------------------------------------------------------------------------------------
\subsection{Small samples and compute}\label{sec:smalln}

\paragraph{The forced variant.} With few points the union bound passes nothing. In round 0, where only the global classes run, no character can pass its test at $d=30$ below $n=17$, the degree-1 floor $2c^2\ln(4Kd/\delta)=16.95$, and no degree-4 character below $n=31$ (\cref{tab:dictionary}). Later rounds have smaller classes and lower floors. In a round where nothing passes, the forced variant adds a guess anyway, at most twice per fit (lines 817--822). If $n\ge48$ and the round's first additive block candidate has $T_B$/threshold above the square of the best character's $z$, it forces that block; both ratios are on a $\chi^2$ scale. Otherwise it adds the character with the largest $z$ among those offered. These are every degree-1 character, the largest pair entry, the largest entry of each row block of degrees 3 and 4, and, from round 1 on, the local subsets above the smallest local threshold (\repo{study/part11/interp/theory/THEORY.md}, \S0). A forced guess passes no test. Forcing a block is on by default (\code{sum_force=True}); the docstring's ``off by default'' (line 73) is out of date.

\paragraph{The paired cross-validation.} At $n\le\code{cv_n}=512$, \code{DFG.fit} chooses between the default and the forced variant by cross-validation (lines 548--571). There are $F=10$ folds at $n\le64$ and 5 above, assigned by a random permutation of the rows. Each variant is fitted on each fold's training rows, without the degree-5/6 search, and predicts the held-out fold. Let $e_{10}$ be the number of points where the default errs and the forced variant does not, and $e_{01}$ the reverse. The forced variant is chosen if $e_{10}-e_{01}>2\sqrt{e_{10}+e_{01}}$: McNemar's test for paired errors \cite{mcnemar1947note}, one-sided at two standard errors. The chosen variant is then fitted on all $n$ rows, and the fold fits' FLOPs count toward the fit. On the 37 targets of \cref{sec:interp} (seed 0), the cross-validation chose the default variant in 204 of the 222 cells with $n\le512$ and the forced variant in 18 (\repo{study/part11/interp/ATLAS.md}, \S2.3) \status{Numerical}. No level is proved for this choice, because the folds' held-out errors are dependent (\cref{sec:guarantees}).

\paragraph{The kernel background.} The output adds $g\cdot x$, with $g=X^\top r/n$ on the final residual. With no feature but the constant, $r=y-w_0$, and the predictor is $\mathrm{sign}(w_0+g\cdot x)$: the centred Hebbian classifier, one step on all $d$ inputs, the counterpart of the networks' near-kernel component. Once the features fit the labels, $r\approx0$ and $g\approx0$. The background has no fitted parameter and costs $2nd$ FLOPs. Least squares, likewise, is kept to $p\le n/2$ features, and the linear block is admitted only if $n\ge2(p+k)$, so the fit never sits at its interpolation peak.

\paragraph{The budget.} The unit is one training run of a toy network at the same $d$, $\mathrm{GD}(d)=394\,PT$ FLOPs for $P=128d+16768$ parameters and $T$ steps: $1.62\times10^{11}$ at $d=30$ (\cref{sec:metrics}). A fit call gets $\mathrm{GD}\cdot\text{scale}$, with scale 1 at $n>512$. At $n\le512$ the final fit gets $\frac12$ and each of the $2F$ fold fits $1/(4F)$, so the shares add up to one run. Within a call:
\begin{itemize}
\item the exhaustive searches run on the first $m$ rows, with $m$ set so that each costs at most a fixed share of the call's budget. The shares are 15\% for degree 4 (round 0 only), 3\% per round for degree 3, 3\% per round for the two-input extensions, and 3\% each for degrees 5 and 6 (round 0 of the final fit only). The search uses $m=\min(n,\lfloor\text{share}\cdot\mathrm{GD}\cdot\text{scale}/c_{\rm row}\rfloor)$ rows, with $c_{\rm row}$ the code's per-row cost estimate (\cref{prop:compute}), and its threshold uses $m$ (\cref{eq:chartest}). A search whose $m$ falls below 32 rows (64 for degrees 5 and 6) is skipped. At $d=30$ only degrees 5 and 6 are subsampled, to 1374 and 147 rows (687 and 73 at $n\le512$). At $d=60$ the degree-4 search sees at most 2212 rows, and at $d=120$ at most 364, with that $d$'s 40000-step run as the unit;
\item a round starts only while the call has spent less than 0.75 of its budget. The other parts are not budgeted: degrees 1 and 2, the one-input extensions, the local class, the eigendecomposition, the sum blocks, least squares, the Newton steps and the background. Only this cap bounds them, so a call can exceed 0.75 by the cost of its last round.
\end{itemize}
\paragraph{The budget in practice} \status{Numerical}. The cap stopped none of the 6786 fit calls recorded by the false-admission and compute analysis; the largest spent 0.471 of its budget (\repo{study/part11/interp/calib/CALIB.md}, \S3.4, confirmed by \repo{study/part11/interp/calib/VERIFY_CALIB.md}). Over the 4730 fits in the result files, the largest uses 0.519 of one network run (\cref{sec:res-compute}).

\begin{table}[t]
\centering{\scriptsize\setlength{\tabcolsep}{3pt}
\begin{minipage}[t]{0.505\linewidth}
\begin{tabularx}{\linewidth}[t]{@{}l@{\hspace{4pt}}l@{\hspace{5pt}}>{\raggedright\arraybackslash}X@{}}
\toprule
\multicolumn{3}{@{}l}{\emph{Tests}}\\
\code{K} & 4 & global degree in round 0\\
\code{K_later} & 3 & global degree, rounds $\ge1$\\
\code{K_hi} & 6 & top degree, round 0\\
\code{c} & 1 & threshold multiplier\\
\code{delta} & 0.1 & level $\delta$ of every test\\
\code{noise} & \code{residual} & tests scaled by $s=\mathrm{rms}(r)$\\
\multicolumn{3}{@{}l}{\emph{Rounds and classes}}\\
\code{rounds} & 8 & rounds at most\\
\code{max_new} & 64 & characters admitted per round\\
\code{max_feats} & 256 & features (also $\le n/2$)\\
\code{max_ext_bases} & 32 & characters extended\\
\code{local_max} & 12 & largest $|V|$, local class\\
\code{blocks} & True & linear and spectral moves\\
\multicolumn{3}{@{}l}{\emph{Small $n$ and budget}}\\
\code{small_n} & True & forced variant and its test\\
\code{cv_n} & 512 & the test runs at $n\le$ this\\
\code{background} & True & kernel background $g\cdot x$\\
\code{budget} & (0.15, 0.03, 0.03) & shares: degree 4; 3; ext.\ by 2\\
\code{budget_hi} & 0.03 & share: degrees 5, 6 each\\
\multicolumn{3}{@{}l}{\emph{Readout}}\\
\code{readout} & \code{ls} & run as \code{logit}\\
\code{ridge} & 0.001 & ridge of least squares ($\times n$)\\
\code{logit_lam} & 0.0001 & L2 penalty, logistic steps\\
\bottomrule
\end{tabularx}
\end{minipage}\hfill
\begin{minipage}[t]{0.47\linewidth}
\begin{tabularx}{\linewidth}[t]{@{}l@{\hspace{4pt}}l@{\hspace{5pt}}>{\raggedright\arraybackslash}X@{}}
\toprule
\multicolumn{3}{@{}l}{\emph{Sum blocks}}\\
\code{sums} & True & sum blocks\\
\code{sum_nmin} & 32 & smallest $n$ for the block search\\
\code{sum_kmax} & 24 & largest block; the $k_{\max}$ of $N_B$\\
\code{sum_mod} & 4 & modulators $\mu$ at most\\
\code{sum_add} & 2 & additive blocks per round\\
\code{sum_gcap} & 2048 & rows for the search's scores\\
\code{sum_mc} & 0.005 & smallest group over $n$ (at least 3)\\
\code{sum_prod} & True & product and split blocks\\
\code{sum_seeds} & (2, 3, 2) & eigenvector, pair, single seeds\\
\code{sum_inner} & True & weighted blocks $t+a\tau$\\
\code{sum_ntau} & 1 & pairs tried as inner terms\\
\code{sum_beam} & 32 & beam width\\
\code{sum_beam_nmax} & 128 & beam only at $n\le$ this\\
\code{sum_force} & True & forced variant may force a block\\
\code{sum_force_nmin} & 48 & smallest $n$ for a forced block\\
\code{sum_force_r2} & 0 & smallest $T_B/n$ of a forced block\\
\bottomrule
\end{tabularx}
\par\medskip
\begin{tabularx}{\linewidth}[t]{@{}l@{\hspace{5pt}}>{\raggedright\arraybackslash}Xr@{}}
\toprule
\multicolumn{2}{@{}l}{\emph{Fixed in the code}} & line\\
15 & logistic Newton steps & 851\\
$n\ge8p$ & when they run & 849\\
$4w$ & where they start & 850\\
0.75 & budget share that ends the rounds & 598\\
10 / 5 & folds at $n\le64$ / above & 553\\
$2\sqrt{e_{10}+e_{01}}$ & margin of the paired test & 567\\
$\frac1{4F}$, $\frac12$ & budget scale: fold fit, final fit & 564, 569\\
32 / 64 & fewest rows: degrees 3, 4 / 5, 6 & 630, 655\\
4; $4d$ & spectral features; smallest $n$ & 707\\
2 & forced guesses per fit & 817\\
$10^{-4}$ & least $\mathrm{mean}(\tilde r^2)/s^2$ for blocks & 734\\
3 & fewest points per group & 721\\
$\pm2,\pm3$ & weights $a$ of inner terms & 273\\
5 & growth steps without gain & 179\\
24 & pairs for inner-term extensions & 180\\
1024 & rows of the beam & 180\\
$3k_{\max}$ & refinement moves at most & 376\\
\bottomrule
\end{tabularx}
\end{minipage}
}
\caption{DFG v1.3's constructor defaults, as run (\code{DFG(readout="logit")}, every other argument at its default), and constants fixed in its code, with their lines \status{Proved, by the code}. Source: \repo{study/part10/hyb_portfolio3/dfg13.py}, read by \repo{paper_dfg/make_tables.py} (\texttt{tab\_params}). The defaults are read from the constructor's signature, and each fixed constant is checked to occur at the line given.}
\label{tab:params}
\end{table}

% ------------------------------------------------------------------------------------------------------------------
\subsection{Simple guarantees}\label{sec:guarantees}

Five statements about DFG can be made exactly. They are P1--P5 of \repo{study/part11/interp/theory/THEORY.md}, restated here for the code as it is. They concern round 0, single tests, the block dictionary and the code's arithmetic; none bounds the test error of a fitted model. Each statement is followed by what it does not cover, and by checks computed by \repo{study/part11/interp/theory/check_theory.py} and stored in \repo{study/part11/interp/theory/check_theory.json} under the keys given. ``By the code'' means a direct reading of the cited lines. Throughout, the rows are iid uniform on $\cube$, as on the toy suites; the column permutation of \code{harness.Permuted} preserves this. We write $L_k=\ln(4K\binom dk/\delta)$ for the log term of the global degree-$k$ class.

\begin{proposition}[Round 0's character and linear tests under the null; P1]\label{prop:null}
\status{Proved}. Let the labels be independent of the rows, with any joint law. In round 0 of any fit call (a final or a fold fit, either variant, any budget scale) the following hold.
\begin{enumerate}[label=(\alph*)]
\item For each degree $k$ the call searches, and every $u>0$,
\[
\Pr\Bigl(\max_{|S|=k}z_S\ge u\Bigr)\le2\tbinom dk\Bigl(\frac\delta{4K\binom dk}\Bigr)^{u^2c^2}.
\]
At $u=c=1$ this is $\delta/(2K)=0.0125$ per degree, for every $d$ and $n$.
\item The linear block's statistic over all $d$ inputs satisfies $\Pr(T_{\rm lin}>d+2\sqrt{dx}+2x)\le e^{-x}$ for every $x>0$. Hence a linear block is admitted with probability at most $\delta+\delta/(2K)$.
\item If round 0 admits nothing, the default variant's model is the constant plus the kernel background. The default variant's model therefore holds a feature other than the constant with probability at most $3\delta/4+\delta=0.175$, plus the probabilities that round 0 admits a spectral feature or a sum block.
\end{enumerate}
\end{proposition}
\begin{proof}
(a) Condition on $y$. Round 0's fit is on the constant alone, so $r=y-w_0$ with $w_0=\sum_ty_t/(n+10^{-9})$, and $r$, $s$ and $\sigma_m$ are functions of $y$. For a fixed nonempty $S$ the signs $\varepsilon_t=\chi_S(x_t)$ are iid uniform and independent of $y$. The statistic is $\frac1m\sum_{t\le m}r_t\varepsilon_t$ and its threshold $c\sigma\sqrt{2L_k/m}$, where $\sigma^2\ge\frac1m\sum_{t\le m}r_t^2$: $\sigma=\sigma_m$ if $m<n$, and $\sigma=s\ge\mathrm{rms}(r)$ if $m=n$. Hoeffding's inequality \cite{hoeffding1963probability}, $\Pr(|\sum_ta_t\varepsilon_t|\ge\lambda)\le2\exp(-\lambda^2/(2\sum_ta_t^2))$, with $a_t=r_t$ and $\lambda=uc\sigma\sqrt{2mL_k}$, gives $\Pr(z_S\ge u\mid y)\le2e^{-u^2c^2L_k}$. A union bound over the $\binom dk$ sets and the expectation over $y$ finish the proof. The degree-5/6 search offers $\frac1m\sum_tr_t\chi_{T_1}(x_t)\chi_{T_2}(x_t)$ only for disjoint $T_1,T_2$ (line 662). That is the same statistic for $S=T_1\cup T_2$, and a set reached by several splits is one event.

(b) Condition on $y$ and let $a=r/|r|$. Since $ns^2\ge|r|^2$, $T_{\rm lin}\le\sum_iZ_i^2$ with $Z_i=\sum_ta_tx_{t,i}$. The $Z_i$ are independent Rademacher sums with $\sum_ta_t^2=1$. For $0\le\lambda<\frac12$ and $\xi\sim N(0,1)$, $\E e^{\lambda Z_i^2}=\E_\xi\prod_t\cosh(\sqrt{2\lambda}\,\xi a_t)\le\E_\xi e^{\lambda\xi^2}=(1-2\lambda)^{-1/2}$. Hence $\ln\E e^{\lambda(\sum_iZ_i^2-d)}\le d\lambda^2/(1-2\lambda)$. This bound on the log-Laplace transform is all that the Chernoff argument in the proof of Lemma~1 of \cite{laurent2000adaptive} uses, and that argument gives the tail. At the code's $x=\ln(1/\delta)$ the bound is $\delta$. The block's inputs are all $d$ unless a degree-1 character passed, which (a) bounds by $\delta/(2K)$.

(c) By the code, a round that admits nothing ends the fit (line 824). In round 0 there are no current characters or blocks, so the extension, local and modulated classes and the products of admitted blocks are empty. Add (a) over degrees 1--6 and (b); the $\delta/(2K)$ term of (b) is (a)'s degree 1.
\end{proof}

\paragraph{Not covered.} The statements are for exact arithmetic; rounding matters only within rounding error of a threshold. Several parts of a fit are not covered.
\begin{itemize}
\item The spectral test: its edge \eqref{eq:spectest} is the asymptotic semicircle edge with a Gaussian correction, not a bound.
\item The sum-block tests: \cref{eq:blocktest} is calibrated for Gaussian noise, but the groups are a data-dependent partition of the rows. The union bound counts every candidate, but the tail of a fixed candidate is not proved to be at most $e^{-L_B}$.
\item The forced variant, which adds features by design.
\item The cross-validation's choice between the variants: a $z$-test at 2 on dependent held-out errors.
\item Rounds $\ge1$: they run only after an admission, and their residual then depends on the rows.
\end{itemize}
\emph{Other input laws} \status{Proved}. The verifier of THEORY.md adds that (a) does not need uniform rows for the searches on all $n$ rows. Round 0's residual sums to zero, up to the $10^{-9}$ term, so $\sum_tr_t\chi_S(x_t)=\sum_tr_t(\chi_S(x_t)-\mu_S)$ with $\mu_S=\E\chi_S$: a sum of independent centred terms of range $2|r_t|$, to which Hoeffding's bound applies unchanged for iid rows from any law (\repo{study/part11/interp/theory/VERIFY_THEORY.md}). This covers the character tests on the dependent one-hot columns of the tabular data (\cref{sec:res-tabular}). It does not cover subsampled searches, whose rows' residuals need not sum to zero, or (b).

\paragraph{Checks} \status{Numerical}.
\begin{itemize}
\item \emph{The character tests.} 16 null configurations of round 0 were simulated, 300 data sets each: at $d=30$ with $n=32$--2048 and three label laws, at $d=15$ with degrees 1--6, and at $d=60$ in a fold fit whose degree-4 search is subsampled. All 280 checks of (a), at levels $u=0.7$--1, held (\texttt{P1.sims}). The bound is loose by about the Gaussian Mills ratio: at $d=30$ the Gaussian approximation gives 0.87--$1.15\times10^{-3}$ per degree, against the bound's $12.5\times10^{-3}$ (\texttt{P1.bound\_over\_gauss\_d30}).
\item \emph{The linear block.} 43 of the 4800 null data sets passed it (0.9\%), against the bound $\delta=10\%$; the $\chi^2_d$ tail at the threshold is 0.92\% at $d=30$ (\texttt{P1.linear\_chi2\_level}).
\item \emph{Full fits on null labels} ($d=30$). With fair labels, 78 of 80 final models are the constant plus the background (\texttt{P1.full\_fits}). One holds a pair at $z=1.06$, the event (a) bounds; the other is a forced-variant model chosen by the cross-validation. With $\Pr(y=1)=0.85$ at $n=128$, 7 of 30 hold spurious spectral or sum-block features (\texttt{P1.nonconstant\_fits}).
\item \emph{The spectral test on biased labels.} Near $n=4d$ its null pass rate grows with $d$. For $\Pr(y=1)=0.8$ it is about 5\%, 47\% and 99\% at $d=30$, 60 and 120 (\repo{study/part11/interp/calib/CALIB.md} with the corrections of \repo{study/part11/interp/calib/VERIFY_CALIB.md}).
\end{itemize}
\Cref{sec:interp-calib} reports DFG's false admissions in full.

\begin{proposition}[The sample size of a parity; P2]\label{prop:parity}
\status{Proved}. Let $y=\chi_S(x)$ with $|S|=k$. Take round 0 of a fit call on $n$ rows in which the degree-$k$ search runs on all of them: $k\le2$ always, and $k=3,4$ when $m_k=n\ge32$. Then $\mathrm{mean}(r\chi_S)=1-w_0\bar y$ and $\mathrm{mean}(r^2)=1-2w_0\bar y+w_0^2$. With $w_0=\bar y$ (dropping the $10^{-9}$), $\chi_S$ passes its test if and only if
\begin{equation}\label{eq:nthr}
n\,(1-\bar y^2)>n_{\rm thr}(k,d)=2c^2\ln\bigl(4K\tbinom dk/\delta\bigr)=2c^2\bigl(k\ln d-\ln k!+\ln(4K/\delta)\bigr)+O(k^2/d).
\end{equation}
Since $\sum_ty_t$ is a sum of $n$ independent fair signs, exactly $\Pr(\text{pass})=\Pr\bigl(|\sum_ty_t|<\sqrt{n(n-n_{\rm thr})}\bigr)$.
\end{proposition}
\begin{proof}
$y\chi_S=\chi_S^2=1$ and $y^2=1$ give the two means. With $w_0=\bar y$, $z_S>1$ reads $1-\bar y^2>c\sqrt{1-\bar y^2}\sqrt{2L_k/n}$. Divide by $\sqrt{1-\bar y^2}$ and square; if $\bar y^2=1$, both sides of the equivalence fail. Finally, $\ln\binom dk=k\ln d-\ln k!+O(k^2/d)$.
\end{proof}
Passing round 0's test is not the same as being in the final model. Admission also needs a place among the round's 64 largest $z$ and room for the feature, which is not proved but never bound in the parity cells. The forced variant can admit $\chi_S$ without passing, and a later round could admit it after another admission. Whether the final model holds $\chi_S$ is therefore the measured statement below. The constant $2c^2\ln(4K/\delta)-2c^2\ln k!$ is not small: $n_{\rm thr}$ is $2k\ln d$ plus 8.8, 6.6 and 3.8 for $k=2$, 3 and 4, so it grows like $2k\ln d$, while $2k\ln d$ alone is 9--44\% below it for $d=15$--120 (\repo{study/part11/interp/theory/VERIFY_THEORY.md}).

\begin{proposition}[A floor for every target; P2$'$]\label{prop:floor}
\status{Proved}. For any labels, in any round, no character passes the test of a class of $N$ characters tested on $m$ rows unless $m>2c^2\ln(4KN/\delta)$. No block passes its test unless $n>\nu+2\sqrt{\nu L_B}+2L_B\ge1+2\sqrt{L_B}+2L_B$.
\end{proposition}
\begin{proof}
By Cauchy--Schwarz, $|\frac1m\sum_{t\le m}r_t\chi_S(x_t)|\le\sigma_m\le\sigma$, so $z_S>1$ needs $c\sqrt{2\ln(4KN/\delta)/m}<1$. For a block, $\sum_vn_v\bar\rho_v^{\,2}\le\sum_t\rho_t^2\le ns^2$: a projection does not increase the norm, $|\rho_t|=|r_t|$, and $s^2\ge\mathrm{mean}(r^2)$. So $T_B\le n$. The same holds for $\tilde r$ in its own units, for products, and for the deflated residual of the second additive search.
\end{proof}
The character half is THEORY.md's P2$'$; the block half is added here. Both bound passing a test, not admission: the forced variant's guesses ($n\le512$) and the linear block's test lie outside these floors. A $k$-parity with $\bar y=0$ attains the character floor, so $n_{\rm thr}(k,d)$ is exactly a parity's sample size. The last column of \cref{tab:dictionary} gives the floors at $d=30$.

\paragraph{Checks} \status{Numerical}.
\begin{itemize}
\item \emph{Every parity cell with $n\ge32$.} \Cref{prop:parity}'s criterion, with each training set's own $\bar y$, predicts the outcome in all 250 parity cells with $n\ge32$: par2--par4 at $d=15$--120, every seed (\texttt{P2.agree}). The test error is 0 where the criterion holds, and 0.488--0.506 elsewhere.
\item \emph{$n=16$.} There $n_{\rm thr}>16$ at every $d$, and nothing passes. The 7 error-free par2 cells at $d\le60$ come from the forced variant, chosen by the cross-validation. With that, the criterion reproduces all 11 measured $\nstar$ (16--56) to within 0.2 (\texttt{P2.pred}).
\item \emph{par4 at $d=30$.} Its $\nstar$ is 52, not 28, because 2 of its 3 seeds have $|\sum_ty_t|>6$ at $n=32$, where $n_{\rm thr}=30.59$. Under \cref{prop:parity}'s exact law that event has probability 0.12.
\end{itemize}

\paragraph{Against the networks} \status{Numerical; Heuristic beyond $d=120$}. $n_{\rm thr}$, which grows like $2k\ln d$, is 19.5--42.0 for $k=2$--4 and $d=15$--120 (\texttt{P2.n\_thr}). DFG's $\nstar$ on the parities moves by at most one grid step over that range (\cref{fig:dscale}). The networks' $\nstar$ grows as $d^{1.03}$, $d^{1.96}$ and $d^{2.40}$ for $k=2$, 3 and 4 (\cref{tab:dscale}), about $d^{k-1}$ for $k\le3$. DFG pays in compute instead: an exhaustive degree-$k$ search costs about $n\binom dk$ FLOPs times a constant (\cref{prop:compute}). Under DFG's fixed budget, at $n>512$, the global search can pass no degree-3 character from $d=630$ on and no degree-4 character from $d=179$ on, because there $m_k<32$ or $m_k\le n_{\rm thr}$ (\cref{prop:floor}; \texttt{P5.d\_parity\_unreachable}). Characters of other targets can still enter through the extension and local classes, which run on all rows. At $d=1000$ the test alone would pass the parities of degree 2--4 at $n\approx56$--57. DFG v1.3 skips its degree-3 and degree-4 searches there ($m_3=17$, $m_4=0$), and an exhaustive degree-4 search on 64 rows would cost 55.9 network runs (\texttt{P2.d1000}, \texttt{P5.exhaustive}).

\begin{proposition}[When a character passes in round 0; P3]\label{prop:nglob}
(a) \status{Proved}. In round 0, for a set $S$ searched on all $n$ rows, $v_S=\hat c_S-w_0\bar\chi_S$ with $\hat c_S=\mathrm{mean}(y\chi_S)$ and $\bar\chi_S=\mathrm{mean}(\chi_S)$, and $s^2=1-2w_0\bar y+w_0^2$. Moreover $z_S>1$ iff $n>2c^2L_{|S|}\,s^2/v_S^2$. For iid uniform rows and nonempty $S$, $\E v_S=\hat f(S)\,\bigl(1-1/(n+10^{-9})\bigr)$.

(b) \status{Heuristic}. Replacing $v_S$ by $\hat f(S)$ and $s^2$ by $1-\hat f(\emptyset)^2$, the round-0 test admits $S$ once
\begin{equation}\label{eq:nglob}
n>n_{\rm glob}(S)=2c^2\ln\bigl(4K\tbinom d{|S|}/\delta\bigr)\,\frac{1-\hat f(\emptyset)^2}{\hat f(S)^2}.
\end{equation}

(c) \status{Proved}. In round $j$, $S$ passes the test of a class of $N_{\rm cls}$ candidates tested on all rows iff $n>n_{\rm cls}=2c^2\ln(4KN_{\rm cls}/\delta)\,s_j^2/v_{S,j}^2$ (on a subsample, $m$ and $\sigma_m$ replace $n$ and $s_j$). Hence
\[
\frac{n_{\rm glob}}{n_{\rm cls}}=\frac{L_{|S|}}{\ln(4KN_{\rm cls}/\delta)}\cdot\frac{1-\hat f(\emptyset)^2}{s_j^2}\cdot\frac{v_{S,j}^2}{\hat f(S)^2}:
\]
a smaller union bound, a smaller residual, and the residual's coefficient on $S$.
\end{proposition}
\begin{proof}
(a) is the algebra of \cref{prop:parity} with $\hat c_S$ in place of 1. For the expectation, $\E[w_0\bar\chi_S]=\frac1{n(n+10^{-9})}\sum_{t,u}\E[y_t\chi_S(x_u)]$. Its terms with $u=t$ are $\hat f(S)$, and those with $u\ne t$ vanish because $\E\chi_S=0$. (c) is the same algebra in round $j$.
\end{proof}

\paragraph{Checks} \status{Numerical}. The checks use the 41 targets of \code{designed} and \code{random}, seed 0, at every $n$ of the grid.
\begin{itemize}
\item \emph{The round-0 decision.} Over 8681 (character, $n$) pairs, $n>n_{\rm glob}$ agrees with DFG's round-0 decision in 8409 (96.9\%; \texttt{P3.test0}). The agreement comes mostly from pairs far from $n_{\rm glob}$: within one doubling of it, the two agree in 773 of 1012 pairs (76\%), and more than two doublings away in 99.9--100\% (\repo{study/part11/interp/theory/VERIFY_THEORY.md}). The pass rate is 6\% one doubling below $n_{\rm glob}$ and 94\% one doubling above.
\item \emph{Entry through the round-0 test.} The 249 characters that first enter through it do so at a median $2^{0.26}\,n_{\rm glob}$ (\texttt{P3.first}).
\item \emph{Entry through later-round classes.} Characters that enter through the local, extension or later global classes do so a median 1.4--3.7 doublings below $n_{\rm glob}$, by class (\texttt{P3.entry\_vs\_nglob}). In \cref{prop:nglob}(c)'s factorization the residual factor is the largest: its class medians are 3.2--19, and per character it ranges from 1.2 to 59.5. The union-bound factor has class medians 1.2--1.7, but single local-class characters gain up to 3.3 from it (\texttt{P3.factors}; per-character ranges from \repo{study/part11/interp/theory/VERIFY_THEORY.md}).
\end{itemize}
\Cref{sec:interp-when,sec:interp-order} use $n_{\rm glob}$ for when, and in what order, DFG admits a target's characters.

\begin{proposition}[Sum blocks; P4]\label{prop:blocks}
\status{Proved}. Let $f(x)=\psi(\beta(x))$, a function of $\beta=\#\{i\in S:s_ix_i=1\}=(t+k)/2$, where $t=\sum_{i\in S}s_ix_i$ is a signed sum of $k$ inputs.
\begin{enumerate}[label=(\alph*)]
\item If every value of $t$ that occurs in training is its own group, the span of the block's indicators and the constant contains $f$ exactly at every $x$ whose value of $t$ occurs in training, with at most $k+1$ weights. (DFG's small ridge shrinks the fitted weights slightly.)
\item Every value is its own group iff every value occurs at least $\max(3,\lfloor0.005n\rfloor)$ times. In the large-$n$ limit this holds iff $k\le7$.
\item A sign readout that uses one block is constant on its groups. It therefore errs on at least $\sum_{\rm groups}\min\bigl(\Pr(f=1,\text{group}),\Pr(f=-1,\text{group})\bigr)$ of the inputs.
\item $\hat f(T)=0$ unless $T\subseteq S$, and $\hat f(T)=(\prod_{i\in T}s_i)\,a_{|T|}$ with $a_j=2^{-k}\sum_\beta\psi(\beta)\mathcal K_j(\beta)$, where $\mathcal K_j(\beta)=\sum_l(-1)^{j-l}\binom jl\binom{k-j}{\beta-l}$. The Fourier weight at degree $j$ is $W_j=\binom kja_j^2$.
\item For a residue class, $\psi(\beta)=\pm(1-2[\beta\equiv\kappa\bmod3])$, and $j\ge1$: $|a_j|\le\frac43\,3^{j/2}\,2^{-k}$, so $W_j\le\frac{16}9\binom kj3^j4^{-k}$.
\item Every starting point of the block search is computed from statistics of one or two inputs, or is an earlier candidate grown from such a starting point. In round 0, $\E\tilde M_{ij}=\hat f(\{i,j\})\bigl(1-1/(n+10^{-9})\bigr)$ for $i\ne j$, and $\E g_i=\hat f(\{i\})\bigl(1-1/(n+10^{-9})\bigr)$.
\end{enumerate}
\end{proposition}
\begin{proof}[Proof sketch]
(a) A function $\eta$ of the group equals $\eta(v_{\max})+\sum_{v\ne v_{\max}}(\eta(v)-\eta(v_{\max}))[\text{group}=v]$ on every occupied group.

(b) By the code (\code{_edges}): scanning the values upward, a group closes once it holds $\max(3,\lfloor0.005n\rfloor)$ points, and a short last group is merged into the one before it. The rarest values, $t=\pm k$, have probability $2^{-k}$, which is at least 0.005 iff $k\le7$.

(c) Any combination of the block's features and the constant is constant on each group.

(d) Inputs outside $S$ are independent of $f$. On $S$, substitute $z_i=s_ix_i$; by symmetry the coefficient depends only on $|T|$, and $\mathcal K_j(\beta)=\sum_{z:\,\#\{z_i=1\}=\beta}z_1\cdots z_j$.

(e) Write $[\beta\equiv\kappa]=\frac13\sum_{m=0}^2\omega^{m(\beta-\kappa)}$ with $\omega=e^{2\pi i/3}$. Since $\beta=\sum_{i\in S}(1+z_i)/2$ with independent $z_i$,
\[
\E\bigl[\omega^{m\beta}z_1\cdots z_j\bigr]=\Bigl(\frac{\omega^m-1}2\Bigr)^j\Bigl(\frac{1+\omega^m}2\Bigr)^{k-j},
\]
which vanishes at $m=0$ for $j\ge1$; the terms $m=1$ and $m=2$ are complex conjugates. Moreover $|\omega-1|/2=\sqrt3/2$ and $|1+\omega|/2=\frac12$, so $|a_j|\le2\cdot\frac23\cdot3^{j/2}2^{-k}$.

(f) The starting points are the eigenvectors and the strongest entries of $\tilde M$; the coefficients $g_i$ and single inputs; inner terms on $\tilde M$'s strongest pair; at $n\le128$, a beam started from two-input blocks; and, in the product and split searches, the round's earlier candidates. A two-input block's groups are functions of $x_i+sx_j$, and these span $\{1,x_i+sx_j,x_ix_j\}$. The expectations follow as in \cref{prop:nglob}(a).
\end{proof}

\paragraph{What this says about the residue classes} \status{Numerical; Heuristic where marked}. Take the three residue-class targets f7\_05, h6\_00 and fd\_10, on $k=14$, 13 and 10 inputs, which DFG alone never learns (\cref{sec:res-suites}).
\begin{itemize}
\item \emph{Fourier weight.} Every one of the $2^k-1$ characters on $S$ is nonzero, while $W_1+W_2=1.6\times10^{-6}$, $5.4\times10^{-6}$ and $2.1\times10^{-4}$. These are exact rationals (\texttt{P4.targets}).
\item \emph{The representation is not what fails.} One block with one group per value represents each target with at most $k+1$ weights, and least squares on such a block reproduces seed 0's training labels at $n=16384$ exactly. DFG's own grouping merges the rarest values, which costs 0.1--0.2\% error (\texttt{P4.grouping}). Part 10 found the true block passing the test from $n=128$ \cite{portfoliolearner2026}.
\item \emph{The search is} (\cref{sec:interp-fail}). By (f), its starting points see these targets only through $a_1$ and $a_2$. \status{Heuristic}: by the BBP threshold for a rank-one perturbation of a random matrix \cite{baik2005phase,benaychgeorges2011eigenvalues}, the eigenvector seed needs $n\approx s^2d/(a_2^2(k-1)^2)$: $1.1\times10^7$, $3.1\times10^6$ and $8.6\times10^4$ rows for the three targets, far above the grid's 16384.
\end{itemize}

\begin{proposition}[Compute; P5]\label{prop:compute}
\status{Proved, by the code}. (a) The global searches of a round are charged, with a multiply-add counting 2 (lines 617--658):
\begin{itemize}
\item $2nd$ FLOPs for degree 1 and $2nd^2+nd$ for degree 2;
\item $(d-2)(2m_3d^2+2m_3d)$ for degree 3 and $\binom{d-2}2(2m_4d^2+3m_4d)$ for degree 4;
\item $m_k(2N_1N_2+N_1k_1+N_2k_2)$ for degrees $k=5,6$, with $N_i=\binom d{k_i}$, $k_1=\lceil k/2\rceil$ and $k_2=\lfloor k/2\rfloor$.
\end{itemize}
The exhaustive degree-3 and degree-4 searches therefore cost about $2nd^3\approx12n\binom d3$ and $nd^4\approx24n\binom d4$. The matrix products also compute redundant entries, so one character costs $12n$ or $24n$ FLOPs, not $2n$.

(b) Each budgeted search runs on $m=\min(n,\lfloor\text{share}\cdot\mathrm{GD}\cdot\text{scale}/c_{\rm row}\rfloor)$ rows, where $c_{\rm row}$ is the code's per-row cost estimate. It is $d(2d^2+2d)$ for degree 3; $\binom d2(2d^2+3d)$ for degree 4; $b(2d^2+3d)$ for the two-input extensions of $b$ bases; and $2N_1N_2+N_1k_1+N_2k_2$ for degrees 5 and 6. These estimates are at least the charged per-row costs, so each budgeted search costs at most its share of the call's budget. The shares of the $2F+1$ calls at $n\le512$ add up to one run, $\frac12+2F\cdot\frac1{4F}=1$. The unbudgeted parts are bounded only by the round cap, up to the cost of the last round.
\end{proposition}

\paragraph{Checks} \status{Numerical}.
\begin{itemize}
\item \emph{An analytic FLOP model} of the ``parity path'' (round 0 admits $\chi_S$ alone, round 1 stops on the sign rule) reproduces the recorded FLOP count to $10^{-9}$ in 228 of the 275 parity cells. It reproduces 4 more once a false linear block is added, which refits confirm; the other 43 cells are at $n\le32$ (\texttt{P5.check}).
\item \emph{The budget bound.} Over $d=5$--399 and the four budget scales, no budgeted search's charged cost exceeds its share; the largest ratio is $1-1.4\times10^{-8}$ (\texttt{P5.max\_share\_used\_over\_share}).
\item \emph{Against one network run.} $\mathrm{GD}(d)$ is linear in $d$, while an exhaustive degree-4 search on 16384 rows costs 0.072, 1.04 and 6.5 network runs at $d=30$, 60 and 120 (\texttt{P5.exhaustive}). The budget's subsampling is what keeps DFG inside one run.
\end{itemize}

\paragraph{What these statements do not give.} None of them bounds the test error of a fitted model. \Cref{prop:null} covers round 0 only. \Cref{prop:parity} is for noiseless parities, and \cref{prop:nglob}(b) is first order. \Cref{prop:blocks} says what a block can represent and what the search's starting points see, not when the search succeeds. \Cref{prop:compute} counts FLOPs, not CPU time (\cref{sec:res-compute}).

% ------------------------------------------------------------------------------------------------------------------
\subsection{How DFG relates to familiar learners}\label{sec:relates}

\status{Heuristic}, except where a statement restates the code or \cref{sec:guarantees}.
\begin{itemize}
\item \emph{Forward stepwise regression and orthogonal matching pursuit.} DFG's loop is forward selection: score candidates by their correlation with the residual, admit, and refit all coefficients jointly by least squares. The joint refit is what distinguishes orthogonal matching pursuit \cite{pati1993orthogonal} from matching pursuit \cite{mallat1993matching}. DFG differs in three ways.
  \begin{itemize}
  \item Its dictionary is never stored as a matrix. It is enumerated class by class, with a fast transform where one exists (the Walsh--Hadamard transform of the local class).
  \item Its entry rule is a significance test, not ``take the best''. If $\chi_S$ is orthogonal on the sample to the current features, it would remove $nv_S^2$ from the residual sum of squares. Then $z_S>1$ iff that reduction, in units of $s^2$, exceeds $2c^2\ln(4KN/\delta)$. With $c=1$ this is the per-variable penalty $2\ln p$ of the risk inflation criterion \cite{foster1994risk}, and the universal threshold $\sqrt{2\ln p}$ of wavelet shrinkage \cite{donoho1994ideal}, with $p=4KN/\delta$: a Bonferroni rule within each class.
  \item A round admits every passing candidate, up to 64, not one at a time.
  \end{itemize}
\item \emph{The low-degree algorithm and sparse Fourier learning.} The low-degree algorithm \cite{linial1993constant} estimates every Fourier coefficient up to degree $k$ from random examples and predicts with the sign of the truncated expansion. \code{lowdeg} (\cref{sec:refs}) is its thresholded version. DFG is a greedy, budgeted, residual-based version of it.
  \begin{itemize}
  \item It estimates coefficients of the residual, not of the labels, and refits the weights instead of using the raw estimates.
  \item It adds classes that grow from what it has: the extensions and the local transform.
  \item It adds a second dictionary, the signed-sum blocks.
  \end{itemize}
  Kushilevitz and Mansour \cite{kushilevitz1993learning}, building on Goldreich and Levin \cite{goldreich1989hardcore}, find every large Fourier coefficient of any degree in time polynomial in $d$, but they need membership queries. DFG has only random examples, so it enumerates each class. A degree-$k$ class costs about $n\binom dk$ FLOPs (\cref{prop:compute}), and characters above degree 6 are reached only through the extensions, the local class and the blocks.
\item \emph{Statistical queries.} Each character test estimates the correlation $\E[r\chi_S]$ of the residual with a fixed function of $x$: a correlational statistical query \cite{bshouty2002using,kearns1998sq}, asked of the residual instead of the label. The union bound is the tolerance budget for all $N$ queries of a class, answered from the same $n$ rows; that is why the sample cost is logarithmic in $N$ (\cref{prop:parity}). For parities, statistical-query learners need a number of queries, or an inverse tolerance, polynomial in the number of candidate parities \cite{blum1994weakly}. DFG pays the query count in compute (the enumeration) and only its logarithm in samples. Analyses of gradient training on sparse parities find costs polynomial in $d$, with an exponent that grows with $k$ \cite{barak2022hidden,abbe2023leap}.
\item \emph{Occam's razor and minimum description length.} The log term $\ln(4KN/\delta)$ is, up to the constant $\ln(4K/\delta)$, the number of nats needed to name one candidate in its class. In those terms the character test reads: the residual sum of squares it removes, over $2s^2$ (the Gaussian code length it saves), exceeds $c^2\ln(4KN/\delta)$. That is a two-part code in the sense of minimum description length \cite{rissanen1978modeling}, without the usual $\frac12\ln n$ for the coefficient. The parity criterion $n>2c^2\ln(4K\binom dk/\delta)$ is the matching Occam sample size \cite{blumer1987occam}: logarithmic in the size of the class. A fitted model's description length is the sum of its features' names; par4's $\chi_{\{1,2,3,4\}}$ costs $\ln(4K\binom{30}4/\delta)=15.29$ nats at $d=30$.
\item \emph{Boosting.} Like gradient boosting \cite{friedman2001greedy}, and L2 boosting in particular \cite{buhlmann2003boosting}, DFG fits the residual stagewise with base learners from a dictionary. Unlike boosting it has no shrinkage and no early stopping: a candidate enters only if it passes its test, and every weight is refitted jointly after each round, as in orthogonal matching pursuit. Its base learners are fixed Boolean features, not fitted trees. And its loop is short. At $n=16384$, 2048 and 512, all the target characters that DFG carries enter in a single round in 32, 30 and 24 of the 33 fits on the 11 Part 5 targets (\repo{study/part11/interp/order/ORDER.md}) \status{Numerical}.
\item \emph{Gradient-trained networks.} Part 6 designed DFG from what Parts 1--5 had seen the study's networks do \cite{gdcircuit2026,discretelearners2026}:
  \begin{itemize}
  \item they grow features greedily, cheapest first;
  \item a feature's cost is set by its correlation with the residual against the noise of the competing directions;
  \item features on inputs already in use are cheaper;
  \item a near-kernel linear component is present from the start;
  \item on every plateau the network computes the sign of a least-squares fit on the features grown so far.
  \end{itemize}
  DFG keeps this picture and replaces the search: enumeration and explicit tests take the place of noisy gradient steps in weight space. Inside DFG's first round at $n=16384$, the order of its test statistics is positively rank-correlated with the order in which the Part 5 networks acquire the same characters, on all 11 targets (Kendall $\tau$ 0.22--0.81, median 0.54; \repo{study/part11/interp/order/ORDER.md}; \cref{sec:interp-order}) \status{Numerical}. The trade is samples for FLOPs. On a $k$-parity DFG needs $n_{\rm thr}(k,d)$ samples, which grow like $2k\ln d$, and about $n\binom dk$ FLOPs of search; the networks need about $d^{k-1}$ samples for $k\le3$, at a cost per step linear in $d$ (\cref{prop:parity,prop:compute}; \cref{sec:res-dscale}).
\end{itemize}

% ------------------------------------------------------------------------------------------------------------------
\subsection{What is discrete}\label{sec:whatdiscrete}

\begin{table}[t]
\centering\footnotesize\renewcommand{\arraystretch}{1.15}
\begin{tabularx}{\linewidth}{>{\raggedright\arraybackslash}p{4.1cm}>{\raggedright\arraybackslash}X>{\raggedright\arraybackslash}p{3.0cm}}
\toprule
object & how it is chosen & continuous?\\
\midrule
characters $\chi_S$ & enumeration in classes; the union-bound test \eqref{eq:chartest} & no\\
sum blocks: $(S,s)$, fold, inner term $a\tau$, modulator, product partner & the block search (eigenvector, pair and single-input seeds, greedy growth, beam, single moves); the $\chi^2$ test \eqref{eq:blocktest} & no\\
a block's groups & contiguous runs of observed values with at least $\max(3,\lfloor0.005n\rfloor)$ points & no\\
the linear block's inputs & one $\chi^2$ test \eqref{eq:lintest} & no\\
the variant (default or forced) & the paired test on held-out errors & no\\
the settings $K$, $c$, $\delta$, budget shares, \dots & constructor defaults, fixed across all data sets (\cref{tab:params}) & no\\
spectral directions $u$ & top eigenvectors of the residual pair matrix & borderline: real directions, computed in closed form\\
the kernel background $g$ & one Hebbian step, $g=X^\top r/n$ & yes, computed in closed form\\
readout weights $w$ & least squares with a small ridge, then 15 Newton steps of logistic regression & yes: a linear readout\\
\bottomrule
\end{tabularx}
\caption{What DFG v1.3 chooses discretely and what it fits continuously \status{Proved, by the code}.}
\label{tab:discrete}
\end{table}

\Cref{tab:discrete} lists what DFG fits.
\begin{itemize}
\item \emph{Every feature but the spectral ones comes from a finite dictionary.} The candidates are enumerated, or proposed by a combinatorial search over signed subsets, and admitted by a test with an explicit threshold, or, in the forced variant, as the best-scoring candidate when nothing passes. A fitted model can therefore be printed as a list of named features (\cref{sec:interp}).
\item \emph{Only the readout is fitted continuously.} Its weights are fitted by least squares and then by a few Newton steps on a convex logistic loss. The background $g$ is a closed-form statistic of the residual, not a fitted parameter. No feature is fitted by gradient descent, and the readout's 15 Newton steps are the only iterative continuous optimization in a fit.
\item \emph{One borderline case.} The spectral directions $u$ are real vectors, computed in closed form as eigenvectors of the residual pair matrix rather than trained. Each feature $(u\cdot x)^2-1=2\sum_{i<j}u_iu_jx_ix_j$ is a fixed combination of degree-2 characters. The post hoc \code{core} run, which has no spectral features, keeps DFG's results (\cref{sec:hypothesis}).
\end{itemize}
